% 贯穿图上的行走、路径、闭行走与环。边旁数字是经过该边的步序。
\begin{tikzpicture}[panelmark/.style={font=\figurefont\fontsize{8}{10}\selectfont,text=FigureInk,anchor=north,inner sep=0pt,outer sep=0pt},
>=Stealth,
  v/.style={circle,draw=FigureStroke,fill=FigureBlueFill,line width=.8pt,minimum size=6.5mm,inner sep=0pt},
  e/.style={draw=FigureGrid,line width=.8pt},
  h/.style={draw=FigureBlue,line width=1.6pt},
  s/.style={fill=white,inner sep=.8pt,text=FigureBlue,font=\figurefont\fontsize{7.5}{9}\selectfont}]
\newcommand{\base}{\node[v] (1) at (-1.0,0) {$1$}; \node[v] (2) at (0,.8) {$2$};
  \node[v] (3) at (0,-.8) {$3$}; \node[v] (4) at (1.0,0) {$4$};}
\begin{scope}[local bounding box=panelrow]
\begin{scope}[local bounding box=panela]
  \base
  \draw[e] (2)--(4) (3)--(4);
  \draw[h] (1)--node[s]{$1,4$}(2); \draw[h] (2)--node[s]{$2$}(3); \draw[h] (3)--node[s]{$3$}(1);

\end{scope}
\begin{scope}[xshift=3.6cm,local bounding box=panelb]
  \base
  \draw[e] (1)--(3) (2)--(3) (3)--(4);
  \draw[h] (1)--node[s]{$1$}(2); \draw[h] (2)--node[s]{$2$}(4);

\end{scope}
\begin{scope}[xshift=7.2cm,local bounding box=panelc]
  \base
  \draw[e] (1)--(2) (1)--(3) (2)--(3) (3)--(4);
  \draw[h] (2)--node[s]{$1,2$}(4);

\end{scope}
\begin{scope}[xshift=10.8cm,local bounding box=paneld]
  \base
  \draw[e] (2)--(4) (3)--(4);
  \draw[h] (1)--node[s]{$1$}(2); \draw[h] (2)--node[s]{$2$}(3); \draw[h] (3)--node[s]{$3$}(1);

\end{scope}
\end{scope}
% Subfigure marks: 8 pt, centred below each panel; common top 3 mm below this row.
\coordinate (panelrowbase) at ($(panelrow.south)+(0,-3mm)$);
\node[panelmark] at (panela.center |- panelrowbase) {(a)};
\node[panelmark] at (panelb.center |- panelrowbase) {(b)};
\node[panelmark] at (panelc.center |- panelrowbase) {(c)};
\node[panelmark] at (paneld.center |- panelrowbase) {(d)};
\end{tikzpicture}
